\section{可观察性与调试实践}
\subsection{日志与监控流水线}
为了让 Agent 的每次运行变得可追踪，OpenHands 会把终端、浏览器、测试、日志都抽象成统一的事件流。每条事件都包含时间戳、调用的工具、上下文快照与预期结果，便于工程师查找出错点。

\begin{itemize}
    \item 终端日志：记录每一个 shell 命令及其返回码，并会与测试结果一同上传。
    \item 浏览器与网络交互：捕捉点击、表单提交与请求响应，便于复现 HTTP 会话与 API 流程。
    \item 文件/代码 diff：自动保存 Agent 修改的 diff 片段、相关文件路径和上下文。
\end{itemize}

\begin{table}[H]
\centering
\begin{tabular}{p{0.25\textwidth} p{0.35\textwidth} p{0.33\textwidth}}
\toprule
层级 & 用途 & 示例 \\
\midrule
命令级 & 捕捉每个 shell 操作与结果 & \texttt{pip install}、\texttt{pytest tests/foo.py} \\
\midrule
上下文级 & 记录使用的上下文版本与检索数据 & 啮合的 Git commit + Issue summary \\
\midrule
结果级 & 归档测试通过/失败、Diff 分析 & Unit test failed with stack trace \\
\bottomrule
\end{tabular}
\caption{OpenHands 观测层级与日志用途。}
\end{table}

\subsection{调试模式与快速回滚}
当 Agent 计划执行高风险命令（如发布、数据库迁移）时，系统会自动切入调试模式：暂停流程、捕捉当前快照，并把命令提醒人类确认。若命令执行失败，平台可以基于日志和 diff 自动回滚，避免污染主分支。

\begin{knowledgebox}{调试建议}
让 Agent 把每一步的“意图”和“预期结果”写入日志，可以加速人工排查并形成训练样本。
\end{knowledgebox}

\begin{warningbox}{滥用日志可能泄露敏感信息}
调试日志中可能包含 API Key、path、stack，必须设置脱敏与访问权限，避免在审计时泄露隐私。
\end{warningbox}

\subsection{本章小结}
\begin{itemize}
    \item 通过统一的事件流与模板化日志，可以快速缩小故障范围。
    \item 调试模式提供了人工确认与自动回滚的双重保障。
\end{itemize}

